\chapter{\ifenglish Conclusions and Discussions\else บทสรุปและข้อเสนอแนะ\fi}

\section{\ifenglish Conclusions\else สรุปผลการดำเนินงาน\fi}
โครงงานนี้ได้นำเสนอกระบวนการพัฒนาระบบการแบ่งส่วนเนื้อเยื่อปกติในช่องปากภายใต้เงื่อนไขที่มีรอยโรคปรากฏ (Lesion-aware Oral Tissue Segmentation) โดยใช้สถาปัตยกรรมโครงข่ายประสาทเทียมเชิงลึก DeepLabV3 ร่วมกับตัวสกัดลักษณะเด่น ResNet-50 ซึ่งจากการสรุปผลการดำเนินงานสามารถแบ่งประเด็นสำคัญได้ดังนี้:

\begin{enumerate}
    \item \textbf{ประสิทธิภาพการเรียนรู้ของตัวแบบ:} ตัวแบบขั้นสุดท้ายแสดงให้เห็นถึงศักยภาพระดับสูงในการระบุพิกัดและขอบเขตของเนื้อเยื่อเป้าหมายทั้ง 8 คลาส โดยได้รับค่า F1-Score เฉลี่ยถึง 0.9038 ในกรณีปกติ และ 0.8740 ในกรณีที่มีรอยโรค ซึ่งถือว่ามีความแม่นยำเพียงพอต่อการนำไปประยุกต์ใช้เป็นเครื่องมือสนับสนุนการวินิจฉัยเบื้องต้น
    \item \textbf{ความสำคัญของคุณภาพชุดข้อมูล (Ablation Analysis):} จากการเปรียบเทียบในลักษณะ Ablation Study ยืนยันว่าการยกระดับคุณภาพของข้อมูลคำสอน (Dataset Refinement) จากชุดข้อมูลดั้งเดิมสู่ชุดข้อมูลฉบับปรับปรุง เป็นขั้นตอนที่ส่งผลบวกต่อประสิทธิภาพมากที่สุด (เพิ่มค่า F1 ขึ้นมากกว่า 10\%) ซึ่งชี้ให้เห็นว่าในงานวิจัยด้านการแพทย์ ความถูกต้องของผลเฉลย (Ground Truth) มีความสำคัญไม่น้อยไปกว่าการเลือกใช้สถาปัตยกรรมตัวแบบ
    \item \textbf{ความคงทนต่อรอยโรค (Lesion-awareness):} ระบบสามารถรักษาความแม่นยำในการแบ่งส่วนเนื้อเยื่อปกติได้แม้ในบริเวณที่ถูกรบกวนด้วยรอยโรคที่มีรูปร่างและสีสันหลากหลาย แสดงให้เห็นว่าโมเดลสามารถเรียนรู้ลักษณะทางกายวิภาค (Anatomical context) ของช่องปากได้อย่างแข็งแกร่ง ไม่ถูกบิดเบือนด้วยเงื่อนไขของพยาธิสภาพที่ปรากฏ
\end{enumerate}

\section{\ifenglish Challenges\else ปัญหาที่พบและแนวทางการแก้ไข\fi}
ในการดำเนินโครงงาน ผู้วิจัยได้เผชิญกับอุปสรรคที่มีความซับซ้อนและส่งผลต่อความแม่นยำของระบบ ดังนี้:
\begin{enumerate}
    \item \textbf{ความแปรปรวนของคุณภาพภาพถ่ายทางการแพทย์:} ภาพถ่ายในช่องปากมีความท้าทายสูงในด้านสภาพแสง (Lighting conditions) แสงสะท้อนจากความชื้น (Specular reflections) และการบดบังของโครงสร้างต่าง ๆ ผู้วิจัยได้บรรเทาปัญหานี้ด้วยเทคนิคการปรับสมดุลภาพ (Normalization) และการเตรียมภาพ (Preprocessing) เพื่อลดความแปรปรวนของสีและแสง
    \item \textbf{ปัญหาความคล้ายคลึงระหว่างคลาส (Inter-class Similarity):} การระบุขอบเขตรอยต่อระหว่าง "เหงือก" และ "กระพุ้งแก้ม" ยังคงมีความคลาดเคลื่อนในบางพิกัด เนื่องจากมีลักษณะทางสายตาที่ใกล้เคียงกันมาก การแก้ไขในอนาคตอาจจำเป็นต้องใช้ข้อมูลเชิงตำแหน่ง (Geometrical context) เข้ามาช่วยพิจารณาเพิ่มขึ้น
    \item \textbf{ความซับซ้อนในการทำ Semantic Labeling:} การทำ Ground Truth ในภาพทางการแพทย์จำเป็นต้องใช้ความเชี่ยวชาญสูงและใช้เวลานาน ซึ่งอาจเกิดความคลาดเคลื่อนจากสายตาของพนักงานผู้ทำการกำกับข้อมูล (Labeler error) ซึ่งผู้วิจัยได้แก้ไขด้วยการทำการตรวจสอบซ้ำ (Double verification) เพื่อลดอัตราความผิดพลาดให้เหลือน้อยที่สุด
    \item \textbf{ข้อจำกัดด้านทรัพยากรในการคำนวณ:} การทดสอบและคัดเลือกจากทั้ง 25 โมเดลสถาปัตยกรรมต้องใช้ทรัพยากรเครื่องประมวลผล (GPU) และเวลาในการฝึกสอนที่สูงมาก ทำให้การปรับจูน Hyperparameters ในวงกว้างทำได้ค่อนข้างจำกัด
    \item \textbf{ข้อจำกัดของกระบวนการ Preprocessing:} การลบวัตถุขนาดเล็กโดยใช้ threshold คงที่ (0.05\%) อาจส่งผลให้วัตถุที่ถูกต้องจริงถูกลบออกไป โดยเฉพาะในกรณีที่มีขนาดเล็ก แนวทางแก้ไขคือการพัฒนาเกณฑ์แบบปรับตัวตามคลาส (class-specific หรือ adaptive thresholding)
    \item \textbf{ข้อจำกัดของกระบวนการ Post-processing:} แม้ว่าการใช้เทคนิค connectivity และ rule-based จะสามารถแก้ไขข้อผิดพลาดเชิงโครงสร้าง (structural errors) เช่น noise และ class overlap ได้ แต่ยังไม่สามารถแก้ไขข้อผิดพลาดเชิงความหมาย (semantic errors) ที่เกิดจากการทำนายของโมเดลได้อย่างสมบูรณ์

\end{enumerate}

\section{\ifenglish%
Suggestions and further improvements
\else%
ข้อเสนอแนะและแนวทางการพัฒนาต่อ
\fi
}
สำหรับการดำเนินงานในขั้นต่อไป เพื่อยกระดับความสามารถของระบบให้มีประสิทธิภาพสูงสุด ผู้วิจัยมีข้อความแนะนำดังนี้:
\begin{enumerate}
    \item \textbf{การประยุกต์ใช้เทคนิค Explainable AI (XAI):} เพื่อสร้างความเชื่อมั่นในการใช้งานทางการแพทย์ ควรมีการนำเทคนิคอย่าง Grad-CAM มาใช้เพื่อแสดงผลให้ผู้ใช้งานเห็นว่าโมเดลใช้ส่วนใดของภาพในการตัดสินใจแบ่งส่วนเนื้อเยื่อ ซึ่งจะช่วยเพิ่มระดับความน่าเชื่อถือ (Trustworthiness) ให้กับทีมแพทย์
    \item \textbf{การผสมผสานตัวแบบ (Ensemble Learning):} การนำข้อดีของโมเดลหลายสถาปัตยกรรมมาทำงานร่วมกัน อาจช่วยลดอัตราความผิดพลาดในคลาสที่มีความยากสูงอย่าง Hard Palate หรือคลาสที่มีตัวอย่างน้อย (Short-tail classes) ได้ดียิ่งขึ้น
    \item \textbf{การพัฒนาระบบคัดกรองอัตโนมัติ (Automated Triage System):} ต่อยอดจากการแบ่งส่วนเนื้อเยื่อไปสู่การระบุประเภทของรอยโรคโดยตรง เพื่อให้ระบบสามารถส่งต่อเคสที่มีความเสี่ยงสูง (High-risk lesions) ให้กับผู้เชี่ยวชาญได้อย่างรวดเร็ว
    \item \textbf{การพัฒนาลบวัตถุขนาดเล็กแบบปรับตัว (Adaptive Thresholding):} การพัฒนาเทคนิคการกำหนดเกณฑ์ขนาดวัตถุแบบปรับตัวตามบริบทของภาพ อาจช่วยลดผลกระทบจากการลบวัตถุที่ถูกต้องจริงออกไป
    \item \textbf{การใช้ข้อจำกัดด้านความข้างเคียงเชิงโครงสร้าง (Structural Adjacency Constraints):} กำหนดเงื่อนไขการตรวจสอบความต่อเนื่องของเนื้อเยื่อ (Connectivity Rules) สำหรับคลาสที่มีความสัมพันธ์เฉพาะ เช่น คลาส C จะต้องมีขอบเขตรอยต่อติดกับคลาส A และ B เท่านั้น หากพบการทำนายคลาส C ในบริเวณที่ไม่มีความเกี่ยวข้องกับสองคลาสข้างต้น ระบบจะทำการประเมินและปรับปรุงคลาสนั้นใหม่ตามบริบทแวดล้อม (Context-aware refinement)
\end{enumerate}
